\documentclass[a4paper,12pt,oneside]{article}
\makeatletter
\def\input@path{{../../}}
\makeatother
\usepackage{discrete-practice}

\begin{document}

\PracticeHeader{1}{10}{Классы булевых функций и полнота системы}

\Task{Задание 1}
Какие из нижеперечисленных функций:
\begin{itemize}
    \item[(а)] являются линейными?
    \item[(б)] являются монотонными?
    \item[(в)] самодвойственные?
    \item[(г)] сохраняют ноль?
    \item[(д)] сохраняют единицу?
\end{itemize}

\begin{enumerate}[label=\bfseries f\arabic*)]
    \item $f_1(a,b,c,d)=(4,5,7,10,11,12,13,14)$
    \item $f_2(a,b,c,d)=(5,7,11,13,14,15)$
    \item $f_3(a,b,c,d)=\overline{a}bc+\overline{c}d+ac\overline{d}$
\end{enumerate}

\noindent
Определить, является ли данный набор функций полным.

\end{document}
